\scanpage{3}
8. \uline{Une estimation (heuristique) du nombre des valeurs propres de frobe-}\\
\uline{nius dans la cohomologie évanescente qui sont des unités $p$-adiques.}

8.1. Soit $X$ un schéma lisse et propre sur le corps parfait $k$, \inserted{de car. $p>0$} et soit $W$ l'anneau des vecteurs de Witt sur $k$. Si on dispose d'une bonne théorie \unsure{de} la cohomologie cristalline $H(X)$ (au sens des catégories dérivées), on \unsure{devra} avoir la relation suivante avec la cohomologie de De Rham
\[
(8.1.1)\qquad H_{\mathrm{DR}}(X)=H(X)\overset{\inserted{L}}{\otimes}_W k
\]
(isomorphisme dans la catégorie dérivée des $k$-modules). On en déduit \unsure{une} formule des coefficients universels: suitesexactes
\diagram{$(8.1.2)$\quad\begin{tikzcd}[column sep=small]
0\arrow[r]&H^i(X)\otimes_W k\arrow[r,dash]&H^i_{\mathrm{DR}}(X)\arrow[r]&\operatorname{Tor}_1^W(H^{i+1}(X),k)\arrow[r]&0
\end{tikzcd}.}
Ce sont des suites exactes fonctorielles en $X$, tout comme l'isomorphisme (8.1.1). Sauf erreur, \deleted{\unsure{tout}} ceci marche dès maintenant pour $p\ne2$, grâce à Berthelot.

\margin{[ On conclut, grâce à \unsure{8.4.3},
\diagram{\begin{tikzcd}[column sep=small]
0\arrow[r]&(H^i(X)\otimes_W k)_{\mathrm{ss}}\arrow[r]&H^i(X,\underline{O}_X)_{\mathrm{ss}}\arrow[r]\arrow[d,phantom,"\text{(nul si }i>n\text{)}"]&({}_p\operatorname{Tors}^{i+1}(X))_{\mathrm{ss}}\arrow[r]\arrow[d,phantom,"\text{nul si }i\geq n"]&0\\
&&\phantom{H}&\phantom{H}&
\end{tikzcd}}
]}
\margin{et
\diagram{\begin{tikzcd}[column sep=small]
0\arrow[r]&(H^{\unsure{i+n}}(X)\otimes_W k)^{\mathrm{ss}}\arrow[r]&H^i(X,\Omega_X^n)^{\mathrm{ss}}\arrow[r]&({}_p\operatorname{Tors}^{i+n+1}(X))^{\mathrm{ss}}\arrow[r]&0
\end{tikzcd}}
où l'indice ${}_{\mathrm{ss}}$ signifie la partie semi-simple pour $F_X=F_X^*\sigma$\\
et \unsure{où} l'exposant ${}^{\mathrm{ss}}$ \rule{6cm}{.3pt} $V_X=\sigma^{-1}F_{X*}$ ]}

Si maintenant $Y$ est une section hyperplane de \deleted{degré assez élevé de} $X$ (supposée maintenant $\subset P_k^r$), on trouve une suite exacte analogue pour $Y$, et un homomorphisme \inserted{$u=(u_0,u_1,u_2)$} de la première dans la seconde. Le diagramme du serpent nous donne alors une suite exacte
\diagram{$(8.1.3)$\quad\begin{tikzcd}[column sep=small,row sep=large]
\inserted{0}\arrow[r]&\inserted{\operatorname{Ker}u_0}\arrow[r]&\inserted{\operatorname{Ker}u_1}\arrow[dll,bend left=15]&&\\
\inserted{\operatorname{Ker}u_2}\arrow[r]&E(Y)\otimes_W k\arrow[r]&E_{\mathrm{DR}}(Y)\arrow[r]\arrow[d,phantom,"\rotatebox{90}{$=$}"]&\operatorname{Coker}u_2\arrow[r]&0\\
&&\inserted{\operatorname{Coker}u_1}&&
\end{tikzcd},}
où on a posé
\[
E(Y)=\operatorname{Coker}(H^{n-1}(X)\xrightarrow{\inserted{u_0}}H^{n-1}(Y)),\qquad E_{\mathrm{DR}}(Y)=\operatorname{Coker}(H^{n-1}_{\mathrm{DR}}(X)\xrightarrow{\inserted{u_1}}H^{n-1}_{\mathrm{DR}}(Y)),
\]
$X$ étant supposée connexe de dimension $n$, enfin $u_2$ désigne l'application
\[
(8.1.4)\qquad u_{\inserted{2}}:\operatorname{Tor}_1^W(H^n(X),k)\longrightarrow\operatorname{Tor}_1^W(H^n(Y),k).
\]

\deleted{\margin{\figure{\begin{tikzpicture}[x=1cm,y=1cm]
\node[draw,ellipse,minimum width=.85cm,minimum height=2cm] (n) at (0,0) {8.2.};
\draw (-.15,-.6)--(-.75,-1.7);
\end{tikzpicture}}}
Noter que $\operatorname{Tor}_1^W(-,k)$ est isomorphe à $\operatorname{Hom}_W(k,-)$ (éléments de $-$ annulés par $pW$), et que \deleted{\unsure{$H^n($}} $\operatorname{Tors}^n(Y)$ ($=$\deleted{\unsure{(}}sous-module de torsion de $Y$) est dual (si on admet la \uline{dualité de Poincaré} habituelle pour la cohomologie cristalline, y compris les formules pour la torsion) de $\operatorname{Tors}^{n-1}(Y)$, et}
\margin{il est maladroit d'utiliser ici la dualité de Poincaré \deleted{\illegible{1}}}

\mspage{2}
